\section{Introduction}
\label{sec:introduction}

Air-quality monitoring networks must forecast pollutant concentrations when a target sensor stops reporting. Consider a station whose PM2.5 and PM10 measurements have been absent for the last day. Its earlier history, local meteorological measurements and contemporaneous readings from other stations may remain available. A pretrained forecaster can consume this incomplete history directly, while a repair procedure can infer the missing target values before prediction. The useful intervention is the one that improves the future forecast under the information available at the forecast origin.

Existing methods provide several relevant capabilities. Neural imputers such as SAITS and TimeMixer++ learn temporal and multivariate representations, and MoTM transfers pretrained imputation representations~\citep{du2023saits,wang2025timemixerpp,motmsoftware}. Imputation quality and downstream prediction quality can differ~\citep{lemorvan2021good}, and downstream evaluation is already included in time-series imputation benchmarks~\citep{du2024tsibench}. Chronos-2 can share information across related series and accept multivariate histories~\citep{ansari2025chronos2}. These capabilities leave a concrete deployment choice: how much of an incomplete history to repair, which observations to preserve, and how to use a statistical forecast alongside the frozen model.

We study a fixed procedure for forecasting after target-sensor outages. Its input includes an observed historical prefix for fitting simple statistical models and a longer, possibly incomplete history ending at the forecast origin. The procedure repairs only recent missing target entries, retains the available long history, and combines the frozen forecaster's point prediction with a statistical forecast. The target is downstream MAE, with MSE reported throughout. The forecasting backbone is unchanged, and current hidden measurements and future outcomes are excluded from the procedure.

Two difficulties motivate this construction. First, unavailable target values must be inferred from a changing set of observed variables. Replacing all variables can also alter useful auxiliary evidence, while restricting the forecaster to a short repaired window discards older observations. Second, a plausible repaired trajectory can still induce an inaccurate forecast. A separate statistical prediction describes another use of the same pre-origin evidence, but combining predictions helps only under suitable component errors. Both effects require matched long-history and statistical-fusion controls.

We define \method{}, which uses prefix-fitted ridge regression on available local and peer observations, inserts the estimated target values into a protected long history, and equally combines the resulting Chronos-2 forecast with a prefix-fitted vector autoregression (VAR). In the motivating PM example, reporting peers estimate the unavailable recent PM values; older target measurements remain in the model context. The VAR branch filters the observed recent history and extrapolates its final state. Both branches end at the same forecast origin, and the fusion weight is fixed across evaluation cases.

The paper makes three concrete contributions. Section~\ref{sec:method} specifies an executable repair-and-fusion procedure with explicit fitting, observation-preservation and information boundaries. Section~\ref{sec:properties} characterizes its preserved inputs and the standard squared-error condition governing fusion. Section~\ref{sec:experiments} evaluates one fixed method across natural outages, artificial target outages and a longer forecasting horizon, with additional component-transfer results on parking availability. On 45 natural-outage windows from 12 Beijing stations, the method obtains MAE/MSE of 0.388053/0.244688, improving on the matched long-native-plus-VAR control by 5.98\%/7.26\%. The other panels delimit this result: different fixed controls are stronger on artificial outages and at horizon 96. The procedure and application focus are selected from development evidence, whose interpretation is made explicit in the experimental design.

\section{Related Work}
\label{sec:related}

\paragraph{Imputation and downstream prediction.}
\citet{lemorvan2021good} analyze the relationship between imputation and prediction with missing covariates. SAITS uses self-attention for imputation~\citep{du2023saits}; TimeMixer++ provides representations for several time-series tasks, including imputation~\citep{wang2025timemixerpp}; and MoTM uses continuous modeling for pretrained time-series imputation~\citep{motmsoftware}. TSI-Bench evaluates imputation methods under multiple missingness conditions and includes downstream tasks~\citep{du2024tsibench}. Our evaluation fixes the forecasting backbone and separately controls the repair region, the forecaster's history and the auxiliary variables it receives.

\paragraph{Missingness-aware and pretrained forecasting.}
Chronos-2 uses group attention to exchange information among related series, variables and covariates~\citep{ansari2025chronos2}. S4M incorporates missingness representations into a trained forecasting architecture~\citep{peng2025s4m}. Our procedure uses an existing Chronos-2 checkpoint and fits only external statistical models. A low-rank adaptation control~\citep{hu2021lora} tests whether updating a small set of additional parameters provides a stronger comparison. Related-series conditioning and adaptation are existing capabilities; neither is treated as an original primitive here.

\paragraph{Statistical repair and forecast combination.}
Ridge regression stabilizes linear estimation with correlated predictors~\citep{hoerl1970ridge}. Linear forecast combination has a long history~\citep{bates1969combination}, and the squared-error decomposition for averaged predictions is established~\citep{krogh1994ensembles}. We use these components to define a specific target-repair procedure for a frozen multivariate forecaster. The empirical question is whether this fixed combination adds value beyond the same long history, the same VAR branch, and strong alternative repairs. We retain those comparisons when their performance is similar to or better than the proposed procedure.
